\documentclass{article}

% Language setting
\usepackage[english]{babel}

% Set page size and margins
\usepackage[letterpaper,top=2cm,bottom=2cm,left=3cm,right=3cm,marginparwidth=1.75cm]{geometry}

% Useful packages
\usepackage[utf8]{inputenc}
\usepackage{amsmath, amssymb}
\usepackage{geometry}
\usepackage{setspace}
\usepackage{lmodern}

\title{Software for Data Science: Julia}
\date{27/10/2025}

\begin{document}
\maketitle

\section*{Basics}

\begin{enumerate}
    \item Assign the number \(22\) to a variable called \texttt{age}.
    
    \item Use the command \texttt{@show} to print the value. 

    \item Use \texttt{typeof()} to check what kind of data is stored in \texttt{age}.

    \item Reassign \texttt{age} to the text \texttt{"Twenty two"}.
    
    \item Again, print the value and the kind of data stored in \texttt{age}.
    
    \item \textit{Question:} What happened ?
\end{enumerate}

\section*{Unicode Variables}

\begin{enumerate}
    \item Define two variables \(\alpha = 7\) and \(\beta = 3\) using unicode.
    
    \item Create a new variable \(\gamma = \alpha + \beta\).
    
    \item Print \(\gamma\) with a message using \texttt{println("γ = \$γ")}.
    
    \item Use \texttt{typeof()} to verify the kind of data of \(\alpha, \beta,\) and \(\gamma\).
\end{enumerate}

\section*{Printing Output}

\begin{enumerate}
    \item Assign your name to a variable \texttt{name} and a number \(year = 2025\).
    
    \item Use \texttt{println()} and \texttt{@show} to output \texttt{name} and \texttt{year}.
    
    \item Compare the outputs produced by \texttt{println} and \texttt{@show}.
\end{enumerate}

\section*{Arithmetic Operations}

\begin{enumerate}
    \item Assign \(a = 8\) and \(b = 3\).
    
    \item Compute and print the following results:
    \[
    a + b, \quad a - b, \quad a \times b, \quad a / b, \quad a // b, \quad a \bmod b
    \]
    
    \item Identify which results represent integer kind of data ($\mathbb{N}$), which represent rational kind ($\mathbb{Q}$), and which represent floating-point kind ($\mathbb{R}$).
    
    \item Evaluate:
    \[
    2 + 3 \times 4^2 \quad \text{and} \quad (2 + 3) \times 4^2
    \]
    Explain how parentheses change the result.
    
    \item Use the modulo operation to check if a number is even or odd:
    \begin{verbatim}
    n = 17
    if n % 2 == 0
        println("Even")
    else
        println("Odd")
    end
    \end{verbatim}
    
    \item Calculate the average of \(x = 10\), \(y = 20\), and \(z = 35\), and print the result with its kind of data.
\end{enumerate}

\section*{Logical Operators}

\begin{enumerate}
    \item Let \(p = \texttt{true}\) and \(q = \texttt{false}\).
    
    \item Compute and display the results of:
    \[
    p \ \&\&\ q, \quad p \ || \ q, \quad \text{and} \quad !p
    \]
    
    \item Write a short program that checks whether a number \(x\) is both positive and even.
\end{enumerate}

\end{document}